\documentclass{article}
\usepackage{amsmath,amssymb}
\usepackage{xurl}
\usepackage[hidelinks]{hyperref}
% Shared commands for the notes in docs/paper-gaps/.
% Every note loads this file via % Shared commands for the notes in docs/paper-gaps/.
% Every note loads this file via % Shared commands for the notes in docs/paper-gaps/.
% Every note loads this file via \input{command}. Project-wide notation
% belongs here, never in per-note redefinitions.

% --- Project configuration ------------------------------------------------
\newcommand{\projectIssueBase}{https://github.com/LionSR/QICLean/issues/}
\newcommand{\projectPrBase}{https://github.com/LionSR/QICLean/pull/}
\newcommand{\projectDocBase}{https://sirui-lu.com/QICLean/docs/}
% ---------------------------------------------------------------------------

\DeclareUnicodeCharacter{2080}{\ifmmode{}_{0}\else\textsubscript{0}\fi}
\DeclareUnicodeCharacter{2081}{\ifmmode{}_{1}\else\textsubscript{1}\fi}
\DeclareUnicodeCharacter{2082}{\ifmmode{}_{2}\else\textsubscript{2}\fi}
\DeclareUnicodeCharacter{2099}{\ifmmode{}_{n}\else\textsubscript{n}\fi}

\newcommand{\Lean}{\textsc{Lean}}

% A formal declaration name, upright sans, breakable at underscores.
\ExplSyntaxOn
\NewDocumentCommand{\leanid}{m}{
  \ifmmode
    \textsf{\detokenize{#1}}
  \else
    \group_begin:
    \urlstyle{sf}
    \tl_set:Nn \l_tmpa_tl {#1}
    \tl_replace_all:Nnn \l_tmpa_tl {\_} {_}
    \exp_args:NV \path \l_tmpa_tl
    \group_end:
  \fi
}
\ExplSyntaxOff

% Traceability links; use in footnotes, not in the running prose.
\newcommand{\ghissue}[1]{\href{\projectIssueBase#1}{GitHub issue \##1}}
\newcommand{\ghpr}[1]{\href{\projectPrBase#1}{PR \##1}}
\newcommand{\doclink}[1]{\href{\projectDocBase#1.html}{\path{#1}}}

% Domain notation, using \providecommand so a note-local refinement stays
% possible.
\providecommand{\tr}{\operatorname{tr}}

% Dirac notation and the identity operator, as in blueprint/src/macros.
\usepackage{dsfont}
\providecommand{\ket}[1]{|#1\rangle}
\providecommand{\bra}[1]{\langle #1|}
\providecommand{\braket}[2]{\langle #1 | #2 \rangle}
\providecommand{\ketbra}[2]{|#1\rangle\!\langle #2|}
\providecommand{\Id}{\mathds{1}}
\providecommand{\one}{\mathds{1}}
\providecommand{\C}{\mathbb{C}}
\providecommand{\MN}[1]{M_{#1}(\C)}
\providecommand{\MD}{M_D(\C)}

% Verdict marker: \gapnote{<kind>}{<status>}. Kind is the policy
% classification (clarification, local-correction, scope-restriction,
% unfaithful, false-source, open-gap); status is open, wip (elimination
% underway), resolved, or historical. Machine-read by the site index;
% prints a verdict line.
\newcommand{\gapnote}[2]{\par\noindent\textbf{Verdict:} #1 (#2).\par}
. Project-wide notation
% belongs here, never in per-note redefinitions.

% --- Project configuration ------------------------------------------------
\newcommand{\projectIssueBase}{https://github.com/LionSR/QICLean/issues/}
\newcommand{\projectPrBase}{https://github.com/LionSR/QICLean/pull/}
\newcommand{\projectDocBase}{https://sirui-lu.com/QICLean/docs/}
% ---------------------------------------------------------------------------

\DeclareUnicodeCharacter{2080}{\ifmmode{}_{0}\else\textsubscript{0}\fi}
\DeclareUnicodeCharacter{2081}{\ifmmode{}_{1}\else\textsubscript{1}\fi}
\DeclareUnicodeCharacter{2082}{\ifmmode{}_{2}\else\textsubscript{2}\fi}
\DeclareUnicodeCharacter{2099}{\ifmmode{}_{n}\else\textsubscript{n}\fi}

\newcommand{\Lean}{\textsc{Lean}}

% A formal declaration name, upright sans, breakable at underscores.
\ExplSyntaxOn
\NewDocumentCommand{\leanid}{m}{
  \ifmmode
    \textsf{\detokenize{#1}}
  \else
    \group_begin:
    \urlstyle{sf}
    \tl_set:Nn \l_tmpa_tl {#1}
    \tl_replace_all:Nnn \l_tmpa_tl {\_} {_}
    \exp_args:NV \path \l_tmpa_tl
    \group_end:
  \fi
}
\ExplSyntaxOff

% Traceability links; use in footnotes, not in the running prose.
\newcommand{\ghissue}[1]{\href{\projectIssueBase#1}{GitHub issue \##1}}
\newcommand{\ghpr}[1]{\href{\projectPrBase#1}{PR \##1}}
\newcommand{\doclink}[1]{\href{\projectDocBase#1.html}{\path{#1}}}

% Domain notation, using \providecommand so a note-local refinement stays
% possible.
\providecommand{\tr}{\operatorname{tr}}

% Dirac notation and the identity operator, as in blueprint/src/macros.
\usepackage{dsfont}
\providecommand{\ket}[1]{|#1\rangle}
\providecommand{\bra}[1]{\langle #1|}
\providecommand{\braket}[2]{\langle #1 | #2 \rangle}
\providecommand{\ketbra}[2]{|#1\rangle\!\langle #2|}
\providecommand{\Id}{\mathds{1}}
\providecommand{\one}{\mathds{1}}
\providecommand{\C}{\mathbb{C}}
\providecommand{\MN}[1]{M_{#1}(\C)}
\providecommand{\MD}{M_D(\C)}

% Verdict marker: \gapnote{<kind>}{<status>}. Kind is the policy
% classification (clarification, local-correction, scope-restriction,
% unfaithful, false-source, open-gap); status is open, wip (elimination
% underway), resolved, or historical. Machine-read by the site index;
% prints a verdict line.
\newcommand{\gapnote}[2]{\par\noindent\textbf{Verdict:} #1 (#2).\par}
. Project-wide notation
% belongs here, never in per-note redefinitions.

% --- Project configuration ------------------------------------------------
\newcommand{\projectIssueBase}{https://github.com/LionSR/QICLean/issues/}
\newcommand{\projectPrBase}{https://github.com/LionSR/QICLean/pull/}
\newcommand{\projectDocBase}{https://sirui-lu.com/QICLean/docs/}
% ---------------------------------------------------------------------------

\DeclareUnicodeCharacter{2080}{\ifmmode{}_{0}\else\textsubscript{0}\fi}
\DeclareUnicodeCharacter{2081}{\ifmmode{}_{1}\else\textsubscript{1}\fi}
\DeclareUnicodeCharacter{2082}{\ifmmode{}_{2}\else\textsubscript{2}\fi}
\DeclareUnicodeCharacter{2099}{\ifmmode{}_{n}\else\textsubscript{n}\fi}

\newcommand{\Lean}{\textsc{Lean}}

% A formal declaration name, upright sans, breakable at underscores.
\ExplSyntaxOn
\NewDocumentCommand{\leanid}{m}{
  \ifmmode
    \textsf{\detokenize{#1}}
  \else
    \group_begin:
    \urlstyle{sf}
    \tl_set:Nn \l_tmpa_tl {#1}
    \tl_replace_all:Nnn \l_tmpa_tl {\_} {_}
    \exp_args:NV \path \l_tmpa_tl
    \group_end:
  \fi
}
\ExplSyntaxOff

% Traceability links; use in footnotes, not in the running prose.
\newcommand{\ghissue}[1]{\href{\projectIssueBase#1}{GitHub issue \##1}}
\newcommand{\ghpr}[1]{\href{\projectPrBase#1}{PR \##1}}
\newcommand{\doclink}[1]{\href{\projectDocBase#1.html}{\path{#1}}}

% Domain notation, using \providecommand so a note-local refinement stays
% possible.
\providecommand{\tr}{\operatorname{tr}}

% Dirac notation and the identity operator, as in blueprint/src/macros.
\usepackage{dsfont}
\providecommand{\ket}[1]{|#1\rangle}
\providecommand{\bra}[1]{\langle #1|}
\providecommand{\braket}[2]{\langle #1 | #2 \rangle}
\providecommand{\ketbra}[2]{|#1\rangle\!\langle #2|}
\providecommand{\Id}{\mathds{1}}
\providecommand{\one}{\mathds{1}}
\providecommand{\C}{\mathbb{C}}
\providecommand{\MN}[1]{M_{#1}(\C)}
\providecommand{\MD}{M_D(\C)}

% Verdict marker: \gapnote{<kind>}{<status>}. Kind is the policy
% classification (clarification, local-correction, scope-restriction,
% unfaithful, false-source, open-gap); status is open, wip (elimination
% underway), resolved, or historical. Machine-read by the site index;
% prints a verdict line.
\newcommand{\gapnote}[2]{\par\noindent\textbf{Verdict:} #1 (#2).\par}

\title{Domains and Zero Dimension in Convex Functional Monotonicity}
\author{}
\date{2026-10-29}
\begin{document}
\maketitle
\gapnote{clarification}{resolved}

\section{The source convention}

Wolf's Theorem~5.16 considers jointly convex real functionals
$F_d:D_d\subseteq\MN{d}^n\to\mathbb R$ in every dimension, invariant under
simultaneous unitary conjugation and tensoring with a common density matrix.
It concludes $F_{d'}((T(A_i))_i)\le F_d(A)$ for every completely positive
trace-preserving $T:\MN{d}\to\MN{d'}$.\footnote{Source:
\cite[Theorem~5.16, Equations~(5.49)--(5.50)]{Wolf2012Quantum}; local source
\path{Notes/WolfNoteTexSource/ch05_schwarz_inequalities.tex}, lines~993--1035.
Tracked in \ghissue{27}.}
The displayed evaluations implicitly require domain membership.

Convexity of $D_d$, closure under unitary conjugation, and closure under
tensoring with a density state justify the enlarged tuples in the proof.
These closures do not by themselves imply that a marginal belongs to the
lower-dimensional domain. The precise domain theorem therefore requires
$(T(A_i))_i\in D_{d'}$ explicitly. Its partial-trace form similarly requires
$(\tr_C A_i)_i\in D_{\dim S}$. These assumptions express well-definedness
and contain no inequality involving $F$. The full-domain specialization has
no membership assumptions. Coordinate independence identifies different
bases of a matrix space of the same dimension.

\section{Positive input dimension}

The conclusion needs $d\ge1$ if zero-dimensional matrix spaces are allowed.
On full domains, set $F_0=0$ and $F_d=1$ for $d\ge1$. These constant
functions are jointly convex and unitarily invariant. Tensoring with a
density state has positive ancilla dimension and preserves whether the
input dimension is zero, so tensor invariance holds as well.
The unique zero linear map $T:M_0\to M_1$ is completely positive and trace
preserving, but the claimed inequality is $1\le0$.

The theorem is thus stated for positive input dimension, consistent with
the density-state interpretation of the source. Output dimension is unrestricted:
a trace-preserving map from positive input dimension cannot have zero output
dimension. The partial-trace theorem permits a zero-dimensional retained
space and requires only a positive-dimensional discarded space. No density
state exists on a zero-dimensional space.

\section{Proof and verification}

The proof uses the actual Weyl average in Equation~(5.50) and joint convexity.
For rectangular maps, a Kraus isometry is extended to a unitary after
adjoining pure states on a common enlarged space. The input factors are
swapped and the output factors reassociated before discarding the Kraus and
input factors. All identities hold for arbitrary complex matrices; no
positive-semidefinite or invertibility assumption on the tuple components
is added. Singular density ancillas are allowed.

The formal results are
\leanid{Matrix.ConvexMatrixFunctional.partialTraceRight_le},
\leanid{Matrix.ConvexMatrixFunctional.map_le}, and
\leanid{Matrix.ConvexMatrixFunctional.map_le_univ}.
The regressions include the constant zero-dimension counterexample and a
three-component trace functional on arbitrary complex tuples.

\bibliographystyle{plain}
\bibliography{../../blueprint/src/references}
\end{document}
